\renewcommand*{\thefootnote}{\arabic{footnote}}
\setcounter{footnote}{0} \setcounter{page}{1}
The discovery of a scalar particle compatible with the Standard Model
(SM) Higgs boson by the ATLAS and CMS collaborations in
2012~\cite{ATLAS:2012yve,CMS:2012qbp} completed the particle content
of the SM and inaugurated a systematic and ever more precise
experimental study of the Higgs sector. In the years since, the LHC
experiments have measured the mass, spin, parity, and couplings of the
Higgs boson in ever finer detail. Combined analyses of the full LHC
Run-2 dataset by ATLAS and CMS~\cite{ATLAS:2022vkf,CMS:2022dwd},
together with more recent updates based on the complete Run-2
sample~\cite{ATLAS:2026comb,CMS:2026nce}, have established the
couplings of the Higgs boson to third-generation fermions and to the
electroweak gauge bosons at the level of a few to ten percent,
found evidence for its decay to muons, and set increasingly stringent limits on
its couplings to light quarks, its total width, possible CP-violating
admixtures, and on Higgs boson pair production, which probes the
trilinear self-coupling directly.  The LHC is now operating in Run 3
at a centre-of-mass energy of $\sqrt{s}=13.6$~TeV, and the
High-Luminosity LHC (HL-LHC) is expected to ultimately deliver up to
$3$--$4\,\mathrm{ab}^{-1}$ per experiment, which will reduce
the statistical uncertainties on many Higgs-boson observables to the
percent level or below~\cite{Cepeda:2019klc}.

This remarkable experimental progress, and the even higher precision
anticipated at the HL-LHC, places correspondingly stringent demands on
the accuracy of the theoretical predictions used to interpret the
data. For several production and decay channels, the precision of
current measurements is already comparable to, or in some cases better
than, the theoretical uncertainties quoted in the most recent
comprehensive reference for Higgs production cross sections, the CERN
Yellow Report~4 (YR4)~\cite{LHCHiggsCrossSectionWorkingGroup:2016ypw},
compiled by the LHC Higgs Cross Section Working Group (LHCHWG) in
2016. Moreover, YR4 does not provide predictions at
$\sqrt{s}=13.6$~TeV, the current LHC operating energy; an interim
interpolation was subsequently provided by the working
group~\cite{Karlberg:2024zxx} as a stop-gap measure pending a full
update. Both of these facts make clear that a systematic update of the
SM predictions for all Higgs-boson production modes, incorporating the
theoretical progress of the past decade, is both timely and necessary
if the LHC and HL-LHC precision programmes are to be fully exploited.

Since the publication of YR4, the theoretical description of every
production mode considered here has advanced. For gluon fusion, the
dominant mode, the inclusive cross section in the heavy-top limit is
now known exactly at N$^3$LO in
QCD~\cite{Anastasiou:2015vya,Anastasiou:2016cez,Mistlberger:2018etf},
rather than through a threshold expansion. It has been supplemented by
exact finite-quark-mass corrections at NNLO, both for the top-only
contribution~\cite{Czakon:2020vql,Czakon:2021yub,Czakon:2024ywb} and
for the top--bottom interference~\cite{Czakon:2023kqm,Czakon:2024ywb},
by mixed QCD--electroweak corrections to the gluon-induced channel
retaining the full dependence on the electroweak boson
masses~\cite{Bonetti:2016brm,Bonetti:2017ovy,Bonetti:2018ukf,Becchetti:2020wof},
and by an approximate N$^4$LO result obtained in the soft-virtual
approximation~\cite{Das:2020adl}. Vector-boson fusion is available at
N$^3$LO in the structure-function approximation~\cite{Dreyer:2016oyx},
and the non-factorisable corrections neglected in that approximation
have since been computed~\cite{Liu:2019tuy,Dreyer:2020urf,Asteriadis:2023nyl}.
Associated production with a $W$ or $Z$ boson has reached N$^3$LO QCD
accuracy for the Drell--Yan-like
contribution~\cite{Baglio:2022wzu}, and the NLO QCD corrections to the
loop-induced $gg\to ZH$ channel have been completed with full
top-quark mass
dependence~\cite{Altenkamp:2012sx,Hasselhuhn:2016rqt,Davies:2020drs,Chen:2020gae,Wang:2021rxu,Alasfar:2021ppe,Chen:2022rua,Degrassi:2022mro,Davies:2025out},
with first NNLO results in the large-$m_t$ limit appearing
recently~\cite{Davies:2025otz}. For $t\bar{t}H$, the NNLO QCD
corrections~\cite{Catani:2022mfv,Devoto:2024nhl} have been matched to
NNLL soft-gluon resummation and combined with the complete NLO
electroweak corrections~\cite{Balsach:2025jcw}. Single-top associated
production remains known at NLO QCD~\cite{Demartin:2015uha,Demartin:2016axk};
NLO electroweak corrections are also available, but only for the three
production channels taken together, since the separation into channels
is no longer well defined once they are
included~\cite{Pagani:2020mov}. Production in association with bottom
quarks has reached N$^3$LO accuracy in the five-flavour
scheme~\cite{Duhr:2019kwi} and has been matched to the four-flavour
scheme at that order using the FONLL method~\cite{Duhr:2020kzd}; the
recommendation presented below is instead based on the
NLO+NNLL\textsubscript{part} matching of
Refs.~\cite{Bonvini:2015pxa,Bonvini:2016fgf}, which allows an explicit
resummation-scale uncertainty to be assigned.

The parton distribution functions (PDFs) entering all of these
predictions have been refined in parallel. The PDF4LHC21
combination~\cite{PDF4LHCWorkingGroup:2022cjn} is the recommended NNLO
baseline and is used throughout this report. Beyond NNLO, approximate
N$^3$LO sets have been produced by the MSHT and NNPDF
collaborations~\cite{McGowan:2022nag,NNPDF:2024nan}, benchmarked
against one another~\cite{Cooper-Sarkar:2024crx} and subsequently
combined~\cite{Cridge:2024icl}. Separately, the effect of QED
corrections to the DGLAP evolution has been quantified for Higgs
production for the first time~\cite{Manohar:2016nzj,Cooper-Sarkar:2025sqw}.
Both developments bear on gluon fusion more than on any other mode,
since the hard-scattering cross section is known one order beyond the
PDFs and since sub-percent changes in the gluon luminosity are no
longer negligible at the precision now reached. They are introduced in
Section~\ref{sec:setup} and their consequences for the recommendation
are worked out in Section~\ref{sec:ggF}.

The LHCHWG is currently preparing a set of updates which will
collectively replace YR4, published as a series of reports in SciPost
Physics Community Reports and referred to in the following as Report~5
(R5). Where numerical predictions are compared, R5 denotes the recommendations
of the present report and YR4 those of its predecessor. Several of these have already appeared, among them dedicated
studies of vector-boson fusion~\cite{Barone:2025jey}, of
$t\bar{t}H$~\cite{Balsach:2025jcw}, of the $gg\to ZH$
channel~\cite{CampilloAveleira:2025rbh}, and of the modelling of
$b\bar{b}H$ production~\cite{10.21468/SciPostPhysCommRep.18}. The
present report is the contribution to R5 concerned with inclusive
production cross sections. It collects the developments listed above
into an updated set of LHCHWG recommendations at
$\sqrt{s}=7,\,8,\,13,\,13.6$, and $14$~TeV and for Higgs-boson masses
between $120$ and $130$~GeV, covering gluon fusion (ggF), vector-boson
fusion (VBF), associated production with a $W$ or $Z$ boson (VH),
associated production with a top-quark pair ($t\bar{t}H$) or a single
top quark ($tH$), and associated production with bottom quarks
($b\bar{b}H$). For $t\bar{t}H$ and $tH$, predictions are provided at
$13$, $13.6$ and $14$~TeV only. For each production mode we document
the ingredients entering the recommendation, define the procedure by
which they are combined, and state the residual theoretical and
parametric uncertainties. We further place the resulting precision in
the context of the most precise current LHC measurements and of
projected HL-LHC
sensitivities~\cite{Cepeda:2019klc,ATLAS:2022hsp}, in order to assess
to what extent the present theoretical uncertainties will remain
sufficient, or become a limiting factor, over the coming decade of LHC
running. Within this scope the report supersedes
YR4~\cite{LHCHiggsCrossSectionWorkingGroup:2016ypw} and the ad interim
recommendations of Ref.~\cite{Karlberg:2024zxx}; branching ratios,
differential distributions and off-shell effects are not addressed
here.

The chapters that follow differ considerably in length, for two
reasons. For most production modes, the calculations entering the
recommendation and the reasoning behind the choices made in combining
them are discussed at length elsewhere in R5 and in the literature it
builds on. Those chapters therefore concentrate on the ingredients
that define the present prediction, on the combination formula and on
the uncertainty prescription, and point the reader onwards for the
supporting material. Gluon fusion is the exception. No such account
exists elsewhere, and the number of separate decisions that had to be
taken is in any case larger. The renormalisation scheme of the
top-quark mass, the finite quark-mass corrections now available at
NNLO, the mismatch between an N$^3$LO matrix element and NNLO PDFs,
and the inclusion of QED evolution each shift either the central value
or the uncertainty at the percent level, and they do not all act in
the same direction. Section~\ref{sec:ggF} consequently treats each
ingredient and its uncertainty in turn, in the manner of YR4, and is
correspondingly longer than the others. This unevenness reflects where
the work was needed rather than the relative importance of the
production modes.

The remainder of this report is organised as
follows. Section~\ref{sec:setup} defines the common input parameters
and conventions used throughout. Sections~\ref{sec:ggF}--\ref{sec:bbH}
present the state of the art and the resulting recommendations for
ggF, VBF, VH, $t\bar t H$/$tH$, and $b\bar b H$ production,
respectively. Section~\ref{sec:Outlook} discusses the resulting
uncertainties in comparison with YR4 and with current and projected
experimental precision, and Section~\ref{sec:conclusions}
concludes. Numerical tables for all production modes, collision
energies and Higgs-boson masses considered are collected in
Appendix~\ref{app:tables}; for gluon fusion these also contain the
information required to convert the recommendation to approximate
N$^3$LO PDFs. The numbers are additionally available at
\url{https://gitlab.cern.ch/LHCHIGGSXS/LHCHXSWG1/crosssections}.
